%!TEX root = ../main.tex
% APÉNDICE A — Instrumento de validación con motociclistas

\chapter{Instrumento de validación con motociclistas}\label{apx:encuesta}

\begin{guia}
Los apéndices contienen material producido por el autor que respalda
el cuerpo pero lo interrumpiría. Acá va el instrumento completo: para
qué se aplicó, cómo se diseñó, a quién llegó y el cuestionario tal como
lo vieron quienes respondieron. Los resultados se analizan en el
capítulo~\ref{cap:resultados}.
\end{guia}

\section{Objetivo}

Evaluar la comprensión del alcance de MotoSafe GBA, la interpretación
de la cobertura y la procedencia de cada capa, la navegación básica y
los problemas de interfaz en usuarios motociclistas reales.

\section{Participantes y metodología}

El instrumento se aplicó a 12 motociclistas reales, en sesiones
individuales moderadas de forma virtual, con un criterio de inclusión
de mayores de 18 años que circulan en motocicleta al menos una vez por
semana en CABA o GBA. El reclutamiento fue por contacto directo, con
muestreo intencional buscando diversidad de zona y frecuencia de uso,
no probabilístico; por esa razón los resultados son indicativos y no
representan a toda la población motociclista de CABA/GBA, como se
declara en el capítulo~\ref{cap:resultados}. Durante las sesiones no se
registraron nombres, domicilios, recorridos personales ni credenciales
de ninguna persona participante.

\section{Consentimiento}

Texto leído a cada participante antes de comenzar la sesión:

\begin{quote}
``La participación es voluntaria. Se evalúa un prototipo académico, no
su capacidad de conducción. Las respuestas se registran sin datos
identificatorios y se analizarán de manera agregada. Puede abandonar la
actividad en cualquier momento.''
\end{quote}

\section{Tareas observadas}

Para cada persona participante se registró, por tarea: completada, con
ayuda o no completada; tiempo aproximado; error observado; y comentario
literal relevante.

\begin{enumerate}
    \item Buscar una zona o avenida conocida y activar una capa de riesgo.
    \item Comparar dos rutas y explicar la cobertura que muestra el sistema para cada una.
    \item Filtrar la prensa por ``Motos'' y abrir una noticia.
    \item Explicar la diferencia entre ``sin datos'' y un nivel bajo de exposición.
\end{enumerate}

\section{Cuestionario}

\begin{enumerate}
    \item Frecuencia de uso de motocicleta (opción única).
    \item Zona habitual de circulación (abierta).
    \item Fuentes que usa actualmente para decidir recorridos (abierta).
    \item La procedencia y la antigüedad de una capa me ayudan a interpretarla (escala 1--5).
    \item Comprendo que el modelo estima gravedad condicionada a que ocurra un siniestro, y que no predice accidentes futuros (escala 1--5, con explicación en palabras propias).
    \item Puedo distinguir ``sin datos'' de un nivel bajo de exposición (escala 1--5).
    \item El mapa, los filtros y los textos son legibles y navegables (escala 1--5).
    \item ¿Qué le generó confusión durante la sesión? (abierta).
    \item ¿Qué funcionalidad modificaría o agregaría? (abierta).
\end{enumerate}

\section{Análisis y uso permitido}

El análisis de este instrumento se limita a presentar conteos
descriptivos y citas anonimizadas, separando los hallazgos de interfaz
de los de interpretación del riesgo y de las propuestas de producto, sin
generalizar a toda la población motociclista a partir de una muestra de
conveniencia de 12 personas. Los resultados correspondientes se
presentan en la sección~\ref{cap:resultados} de este informe.
